\section{格点数据中的演化}

\subsection{总体特征}

文献 \cite{Orginos:2017kos} 报告了对核子内价部分子分布 $u_v-d_v$ 的约化伪 ITD ${\mathfrak M} (\nu, z_3^2)$ 的探索性格点研究。该文一个惊人的观察是：当作为 $\nu$ 的函数作图时，实部与虚部的数据都靠近各自的普适曲线。数据显示在大 $z_3$ 处没有多项式型的 $z_3$ 依赖。鉴于所探范围内 $z_3^2/a^2$ 从 1 变到约 200，我们把这一结果解释为约化伪 ITD 中{\it 完全不存在}高扭度项。

如文献 \cite{Radyushkin:2017cyf,Orginos:2017kos} 及引言所述，这样的结果对应 Ioffe 时间分布软部 ${\cal  M} (\nu, z_3^2)= M (\nu) {\cal  M} (0, z_3^2)$ 的 $\nu$ 依赖与 $z_3^2$ 依赖的因子化。用 TMD ${\cal F} (x, k_\perp^2)$ 的语言说，这对应其在软 $k_\perp$ 区域内 $x$ 依赖与 $k_\perp^2$ 依赖的因子化。然而，如文献 \cite{Orginos:2017kos} 所观察到的，在小 $z_3$ 值（即 $z_3 \lesssim 6a$）处存在相当明显的 $z_3$ 依赖，可由微扰演化解释。

让我们先考虑实部。它对应函数 $q_v(x)$（价组合，即夸克与反夸克分布之差 $q_v(x) = q(x) -\bar q(x)$；在我们的情形 $q= u-d$）的余弦傅里叶变换
  \begin{align}
   {\mathfrak  R} (\nu)  \equiv
{\rm Re} \,  {\mathfrak  M} (\nu)=  & \int_0^1 dx \
\cos (\nu x)  \, q_v  (x)
\label{MC}
     \end{align}
文献 \cite{Orginos:2017kos} 发现实部数据非常接近由函数
     \begin{align}
    f(x) =  \frac{315}{32} \sqrt{x}  (1-x)^{3}\,  .
    \label{qV}
    \end{align}
生成的曲线 ${\mathfrak R}_f ( \nu)$（见图~\ref{realc}）。这一形状是通过构造归一化 $x^a(1-x)^b$ 型函数的余弦傅里叶变换并拟合数据确定参数 $a,b$ 而得到的。

   \begin{figure}[tbp]
  \centerline{\includegraphics[width=3.3in]{images/data713.pdf}}
     \caption{${\mathfrak  M} (\nu, z_3^2)$ 的实部，$z_3$ 从 $7a$ 到 $13a$。
    \label{M713}}
    \end{figure}

\begin{figure}[tbp]
  \centerline{\includegraphics[width=3.3in]{images/data16.pdf}}
     \caption{${\mathfrak  M} (\nu, z_3^2)$ 的实部，$z_3$ 从 $a$ 到 $6a$。
     \label{M16}}
    \end{figure}

\begin{figure}[tbp]
    \centerline{\includegraphics[width=3.3in]{images/datAll}}
     \caption{${\mathfrak  M} (\nu, z_3^2)$ 的实部作为 $\nu =Pz_3$ 的函数，
     并与方程 (\relax {\ref{MC}})、(\relax {\ref{qV}}) 给出的曲线比较。
        \label{realc}}
    \end{figure}

尽管拟合使用了所有数据点，曲线形状显然由 ${\rm Re}$ ${\mathfrak  M} (\nu, z_3^2)$ 较小的点主导。为给出更细致的说明，我们在图~\ref{M713} 中展示 $7a \leq z_3 \leq 13a$ 范围内的点。可以看到，在有限体积效应变得重要的区域 $\nu \gtrsim 10$ 内，最大 $\nu$ 处的点有些散布。除此之外，几乎所有的点都落在基于 $f(x)$ 的普适曲线上。从这个意义上说，大 $z_3$ 数据中看不到 $z_3$ 演化。

图~\ref{M16} 展示 \mbox{$a \leq z_3 \leq 6a$} 区域的点（注意，格点上 $z_3=0$ 意味着 $\nu=0$，且按定义 ${\mathfrak  M} (0, 0)=1$）。此时所有的点都高于普适曲线。我们记得，微扰演化使伪 ITD 的实部随 $z_3$ 减小而增大。于是可以推测，较小 $z_3$ 点处观察到的较高 ${\mathfrak R}$ 值可能是演化的结果。

格点数据 $z_3$ 依赖的一个典型模式示于\mbox{图~\ref{nu23}}，取``魔数''Ioffe 时间值 $\nu =3\pi/4$——它可由文献 \cite{Orginos:2017kos} 所用的五种不同的 $z_3$ 与 $P$ 组合得到。目测拟合线的形状由不完全伽马函数 $\Gamma (0,z_3^2/30a^2)$ 给出。该函数完全符合如下预期：$z_3$ 依赖在小 $z_3$ 处具有``微扰''的对数 $\ln (1/z_3^2)$ 行为，而在 $z_3$ 大于 $6a$ 后迅速消失。

正如预期，${\mathfrak R} ( \nu, z_3^2)$ 随 $z_3$ 增大而减小。我们还看到演化在大 $z_3$ 处``停止''。在此语境下，基于式 (\ref{qV}) 的整体曲线对应``低标度点''，即微扰演化缺席的区域。

\begin{figure}[tbp]
  \centerline{\includegraphics[width=3.3in]{images/nu3pi.pdf}}
     \caption{$\nu=3\pi/4\approx2.3562$ 时对 $z_3$ 的依赖。
    \label{nu23}}
    \end{figure}

\subsection{构建 $\overline{\rm MS}$ ITD}

于是我们看到，\mbox{图~\ref{nu23}} 的数据在小 $z_3$ 区域显示出对数演化行为。不过，对 $z_3\gtrsim 5a$，\mbox{$z_3$ 行为}开始明显偏离纯对数的 $\ln z_3^2$ 图样。这就给``对数区域''设定了边界 $z_3 \leq 4a$。那么就让我们尝试在该区域使用式 (\ref{MNL}) 来构建 $\overline{\rm MS}$ ITD。

把式 (\ref{MNL}) 中各贡献拆开是有启发性的，记 ${\rm Re} \, {\cal I}  (\nu, \mu^2)  \equiv  {\cal I}_R  (\nu, \mu^2) $。第一部分``演化''部分由
  \begin{align}
 {\cal I}^{\rm ev} _R  (\nu, \mu^2)    =&  {\mathfrak  R} ( \nu, z_3^2) +  \frac{\alpha_s}{2\pi} \, C_F\
 \int_0^1  dw \,   {\mathfrak  R}  (w \nu,z_3^2)   \nn &
 \times    B(w) \,   \ln \left (z_3^2\mu^2 \frac{  e^{2\gamma_E}}{4} \right )
\
\label{MNLev}
 \end{align}
给出（回顾 $ {\mathfrak  R} ( \nu, z_3^2) \equiv {\rm Re} \,  {\mathfrak  M} ( \nu, z_3^2$)），对应文献 \cite{Orginos:2017kos} 使用的领头对数近似。当 $z_3 =2 e^{-\gamma_E}/\mu$ 时对数为零，我们有
   \begin{align}
 {\cal I}^{\rm ev}_R   (\nu, \mu^2)    =&  {\mathfrak  R} ( \nu, (2 e^{-\gamma_E}/\mu)^2) =
  {\mathfrak  R} ( \nu, (1.12/\mu)^2) \  .
\label{IMev}
 \end{align}

当然，这种情况只在如下条件下发生：适当选取 $\alpha_s$ 后，单圈修正的 \mbox{$\ln z_3^2$ 依赖}恰好抵消数据的实际 $z_3^2$ 依赖（在图~\ref{M16} 数据点的散布中可见）。文献 \cite{Orginos:2017kos} 发现这在 $\alpha_s/\pi \approx 0.1$ 时发生。因此，式 (\ref{MNL}) 只在数据对 $z_3$ 显示{\it 对数}依赖的区域（在我们情形即 $z_3 \leq 4a$）才是准确的。

\begin{figure}[tbp]
  \centerline{\includegraphics[width=3.4in]{images/bulamulaN.pdf}}
     \caption{方程 (\ref{MNLo}) 的函数 $B \otimes  {\mathfrak R}_f$（上线）与 $L \otimes  {\mathfrak R}_f$（下线）。
    \label{BMLM}}
    \end{figure}

由于 ${\mathfrak  R}  (w \nu,z_3^2)$ 与 ${\mathfrak R}_f (w \nu)$ 之差是 ${\cal O} (\alpha_s)$，可以在式 (\ref{MNLev}) 中把 ${\mathfrak  R}  (w \nu,z_3^2)$ 替换为 ${\mathfrak R}_f (w\nu)$（回顾 ${\mathfrak R}_f (\nu)$ 对应式 (\ref{qV}) 的 PDF）。${\cal I}  (\nu, \mu^2) $ 的剩余部分（其中已代入 ${\mathfrak R}_f (w \nu)$）
    \begin{align}
{\cal I}^{\rm NL}_R   (\nu)    = & \frac{\alpha_s}{2\pi} \, C_F\
 \int_0^1  dw \,   {\mathfrak R}_f (w \nu)   \nn &
 \times  \left  \{   B(w) \
+   \left [4\,   \frac{\ln (1-w)}{1-w} -2 (1-w)  \right ]_+ \right  \}
 \nn & \equiv  \frac{\alpha_s}{2\pi} \, C_F\, \left  [ B \otimes  {\mathfrak R}_f + L \otimes  {\mathfrak R}_f\right ]
\
\label{MNLo}
 \end{align}
来自超出领头对数近似的修正。

正如我们所讨论的，$L \otimes  {\mathfrak R}_f$ 项反映顶点图演化部分的实际标度小于 $z_3$ 这一事实。为说明其影响，我们在图~\ref{BMLM} 中展示函数 $B \otimes  {\mathfrak R}_f$ 与 $L \otimes  {\mathfrak R}_f$。可以看到后者为负且相当大，其 $\nu$ 依赖与 $B \otimes  {\mathfrak R}_f$ 相似。事实上，在 $\nu <5$ 区域有 $L \otimes  {\mathfrak R}_f \approx -3.5 B \otimes  {\mathfrak  R}_v$。于是这两项的合并效果接近 $-2.5 B \otimes  {\mathfrak R}_f$。结果是，纳入这些项可以近似地处理为带修改重标度因子的 LLA 演化。具体地，我们可以写
    \begin{align}
 {\cal I}_R   (\nu, \mu^2)    \approx &  \, {\mathfrak  R} ( \nu, (2 e^{1.25-\gamma_E}/\mu)^2)
   \approx   \, {\mathfrak  R} ( \nu, (4/\mu)^2)   \  .
\label{IMevL}
 \end{align}
这样，重标度因子相比原来的 LLA 值改变了 4 倍！

我们可以把 $\mu \approx 4/z_3$ 用作指引，但实际的数值计算当然应使用``精确''的式 (\ref{MNL}) 进行。为此，我们选取 $\mu=1/a$，按文献 \cite{Orginos:2017kos} 采用的格距 $0.093$ fm 约为 \mbox{2.15 GeV。}估计式 (\ref{IMevL}) 告诉我们，此标度下的 ITD $ {\cal I}_R   (\nu, \mu^2)$ 应接近 \mbox{$z_3\approx 4a$}（处于 $z_3 \leq 4a$ 区域边界上的距离）处的伪 ITD ${\mathfrak  R} ( \nu, z_3^2) $。取文献 \cite{Orginos:2017kos} 使用的值 $\alpha_s/\pi =0.1$，把完整单圈关系 (\ref{MNL}) 应用于 $z_3 \leq 4a$ 的数据，便生成 ${\cal I}_R  (\nu, (1/a)^2) $ 的各点。

\begin{figure}[tbp]
  \centerline{\includegraphics[width=3.35in]{images/Re035.pdf}}
     \caption{用 $z_3$ 从 $a$ 到 $4a$ 的数据计算的 $\mu = 1/a$ 处的函数 ${\cal I}_R  (\nu, \mu^2) $。
        上曲线对应 CJ15 全局拟合 PDF 的 ITD。
    \label{Msbar16}}
    \end{figure}

从\mbox{图~\ref{Msbar16}}可见，所有各点都以相当小的散布靠近某条普适曲线。曲线本身是把各点用归一化 $N x^a (1-x)^b$ 分布的余弦变换拟合得到的，给出 $a=0.35$ 和 $b=3$。散布的大小说明了 $\nu \leq 4$ 区域内 ITD 拟合的误差。在图~\ref{Msbar16} 中，我们把 $\mu=1/a$ 的 ITD 与由对应 CJ15 \cite{Accardi:2016qay} 全局拟合的全局拟合 PDF 得到的 ITD 作了比较。可以看到我们的 ITD 系统性地低于基于全局拟合 PDF 的曲线。

偏差的原因可以从图~\ref{fabx}理解：图中我们把归一化分布 \mbox{$N x^{0.35} (1-x)^3\equiv q_v (x, \mu=2.15$ GeV)} 与取自标度 $\mu=2.15$ GeV 的 CJ15 \cite{Accardi:2016qay} 和 MMHT 2014 \cite{Harland-Lang:2014zoa} 全局拟合 PDF 比较。不同于 $\sim x^{0.35} $ 函数，这些 PDF 在小 $x$ 处是奇异的，这导致 ITD 在大、中 $\nu$ 值处增强。

拟合 ${\cal I}_R  (\nu, \mu^2) $ 各点时，我们采用了与文献 \cite{Orginos:2017kos} 相同的最简 $N x^a (1-x)^b$ PDF 参数化。原则上可以使用更复杂的 PDF 模型，并在 $\nu \leq 4$ 区域得到实际相同的 ITD 拟合曲线，而 PDF $q_v (x)$ 曲线则略有不同。原因很简单：只有当完整知道整个 $0\leq \nu <\infty$ 区域的 ITD 时，逆余弦傅里叶变换才是唯一的。从有限的 $\nu \leq 4$ 区域作此变换，需要对 ITD 在该区域外的行为或 PDF $q_v (x)$ 的函数形式补充一些假设。我们通过取 $q_v(x)\sim x^a (1-x)^b$ 固定了选择。若采用更复杂的形式（特别是全局拟合 \cite{Accardi:2016qay,Harland-Lang:2014zoa} 所用的形式），$q_v (x)$ 形状如何变化是一个有趣的问题，但超出了本文的范围。

\begin{figure}[tbp]
  \centerline{\includegraphics[width=3.6in]{images/fx035.pdf}}
     \caption{由图~\ref{Msbar16} 数据构建的 $\mu=2.15$ GeV 下 $u_v(x)-d_v(x)$ 曲线，
 并与 CJ15 和 MMHT 全局拟合比较。
    \label{fabx}}
    \end{figure}

与文献 \cite{Orginos:2017kos} 的 LLA 结果比较，我们看到式 (\ref{MNL}) 中大的负单圈修正明显改变了提取出的 PDF，使其离全局拟合 PDF 更远。主要原因是文献 \cite{Orginos:2017kos} 构建的 $z_0=2a$ 伪 ITD 在该文中被当作对应 $\mu\approx 1$ GeV 标度，而按照修改后的重标度关系 (\ref{IMevL})，它应对应 $\mu \approx 4$ GeV。因此要得到 $\mu \approx 2$ GeV 的曲线，需要把它向低 $\mu$ 演化。

不过，文献 \cite{Orginos:2017kos} 的指导性思想——$\overline{\rm MS}$ ITD ${\cal I}_R(\nu,\mu^2)$ 可以通过适当重标度 $\mu = k/z_3$ 从约化伪 ITD ${\mathfrak R} (\nu,z_3^2)$ 得到——在取 $k\approx 4$ 时对所有 $z_3 \leq 6a$ 都以相当好的精度成立。按这一重标度关系，$\mu  =1/2a  \approx 1$ GeV 的 ITD 对应 $z_3 \approx  8a$ 的约化伪 ITD。正如我们所讨论的，演化停止的边界点约为 $z_3 \approx 6a$。因此该距离处的伪 ITD 由对应\mbox{式 (\ref{qV})}普适拟合函数 $f(x)$ 的 ITD ${\mathfrak R}_f (\nu)$ 给出。这一结果也可由基于\mbox{式 (\ref{MNL})}的直接数值计算得到。

利用式 (\ref{MNL}) 还可以把 $\overline{\rm MS}$ ITD 向 $\mu=1/2a$ 以下演化，所得函数将随 $\mu$ 变化。另一方面，伪 ITD 在 $z_3 \gtrsim 6a$ 时不随 $z_3$ 改变。因此在该区域内，重标度联系 ${\cal I}_R (\nu, \mu^2) \approx {\mathfrak R} (\nu, (4/\mu)^2)$ 随 $\mu$ 减小而越来越不准确，最终失去意义。

\subsection{虚部}

\begin{figure}[tbp]
  \centerline{\includegraphics[width=3.4in]{images/im713.pdf}}
     \caption{${\mathfrak  M} (\nu, z_3^2)$ 的虚部，$z_3$ 从 $7a$ 到 $13a$。
    曲线对应 $q (x)  +  \bar q (x)= f(x) + 2 \bar q (x)$，其中 $f(x)$ 由式 (\ref{qV}) 给出，
    $\bar q (x)$ 由式 (\ref{barq}) 给出。
    \label{imor713}}
    \end{figure}

伪 ITD 的虚部可以用类似方式考虑。它对应夸克与反夸克分布之和 $q(x) + \bar q (x) $ 所给函数的正弦傅里叶变换
  \begin{align}
   % {\mathfrak  J} (\nu)  \equiv
{\rm Im} \,  {\mathfrak  M} (\nu)=  & \int_0^1 dx \
\sin  (\nu x)  \, [q (x)  +  \bar q (x)]
\label{MS}
     \end{align}
该函数与价组合 $q_v(x) = q(x) -\bar q(x)$ 相差 $2 \bar q (x) = 2[ \bar u (x) - \bar d(x)]$。图~\ref{imor713} 展示大 $z_3$ 值 $z_3 \geq 7a$ 的数据。与实部情形一样（见图~\ref{M713}），$\nu \lesssim 10$ 的点靠近一条普适曲线。写出 $q (x)  +  \bar q (x) =q_v(x) +2 \bar q (x)$ 并取式 (\ref{qV}) 的 $f(x)$ 作为 $q_v (x)$，我们发现
  \begin{align}
    \bar q (x) \approx  0.1 \,[20  x \, (1-x)^{3}]\,  .
    \label{barq}
    \end{align}
注意在文献 \cite{Orginos:2017kos} 中拟合是对全部 $z_3$ 数据点做的（即也纳入了 $z_3 \leq 6a$ 的点），$\bar q (x) $ 的整体系数得到 0.07 而非 0.1。

图~\ref{imor06} 展示 $z_3 \leq 4a$ 的数据。可以看到所有这些点都低于由 $z_3 \geq 7a$ 数据拟合得到的曲线。这与下述事实一致：在 $\nu \lesssim 6$ 区域，微扰演化使伪 ITD 的虚部随 $z_3$ 减小而减小。注意单圈关系对整个函数 ${\mathfrak M}= {\rm Re } \, {\mathfrak M} + i \, {\rm Im } \, {\mathfrak M}$ 成立，所以只需分离其实部与虚部即可，$\overline {\rm MS}$ 函数 ${\rm Im}\,  {\cal I} (\nu, \mu^2) \equiv {\cal I}_I (\nu, \mu^2)$ 的构建方式与实部相同。

结果示于\mbox{图~\ref{imev}}。同样，所有各点以相当小的散布靠近一条普适曲线。所示曲线对应价分布 \mbox{$q_v (x, \mu =1/a) =N x^{0.35} (1-x)^3$}（由实部研究得到）与反夸克贡献 $2 \bar q (x, \mu =1/a$）之和的正弦傅里叶变换。后者由拟合得出为 $\bar q (x, \mu =1/a=2.15$ GeV) $= 0.07 [20 x (1-x)^3]$。

\begin{figure}[tbp]
  \centerline{\includegraphics[width=3.5in]{images/imz06.pdf}}
     \caption{${\mathfrak  M} (\nu, z_3^2)$ 的虚部，$z_3$ 从 $a$ 到 $6a$。
曲线同图~\ref{imor713}。
    \label{imor06}}
    \end{figure}

\begin{figure}[tbp]
  \centerline{\includegraphics[width=3.4in]{images/imevoS.pdf}}
     \caption{用 $z_3$ 从 $a$ 到 $4a$ 的数据计算的 $\mu = 1/a$ 处的函数 ${\cal I}_I  (\nu, \mu^2) $。
    曲线见正文描述。
    \label{imev}}
    \end{figure}
